\PassOptionsToPackage{quiet}{xeCJK}
\documentclass[answer]{THUExam}
\usepackage{HTNotes-math}
\usepackage{pifont}
\usepackage{ulem}

\usepackage{tikz}
\usetikzlibrary{positioning, calc, shapes.geometric, shapes, shapes.multipart, arrows.meta, arrows, decorations.markings, external, trees}

% =================================================
%       PDF信息
% =================================================
\title{高等数学A (二)}
\author{}
\Subject{作业}
\Keywords{高等数学, 多元函数微分学, 几何应用, 极值}

% =================================================
%       试卷头信息
% =================================================
\Year{2025--2026}
\Semester{春季}
\Course{高等数学A}
\Time{}
\Type{A}

\begin{document}
\vspace{1cm}
\makehead

\vspace{1cm}

% =================================================
%       选择题：几何应用与海瑟矩阵
% =================================================
\makepart{选择题}{每小题~4~分，共~28~分}

\begin{problem}
曲线 $\Gamma: x=\sin t,\ y=\cos 2t,\ z=\mathrm{e}^t$ 在 $t=0$ 处的切向量为\pickout{A}
\options
{$(1,0,1)$}
{$(0,1,1)$}
{$(1,-2,1)$}
{$(1,0,\mathrm{e})$}
\end{problem}
\begin{note}
切向量为 $(x'(t),y'(t),z'(t))=(\cos t,-2\sin 2t,\mathrm{e}^t)$，代入 $t=0$ 得 $(1,0,1)$。
\end{note}
\bigskip

\begin{problem}
曲线 $\Gamma: x=t,\ y=\ln(1+t),\ z=\dfrac{1}{1+t}$ 在 $t=0$ 处的法平面方程为\pickout{B}
\options
{$x+y+z=1$}
{$x+y-z=-1$}
{$x-y+z=1$}
{$x+y-z=1$}
\end{problem}
\begin{note}
切向量 $s=\bigl(1,\frac{1}{1+t},-\frac{1}{(1+t)^2}\bigr)\big|_{t=0}=(1,1,-1)$，对应点为 $(0,0,1)$。
法平面：$1\cdot(x-0)+1\cdot(y-0)-1\cdot(z-1)=0$，整理得 $x+y-z=-1$。
\end{note}
\bigskip

\begin{problem}
曲面 $z=\mathrm{e}^{x}\sin y$ 在点 $\bigl(0,\frac{\pi}{2},1\bigr)$ 处的切平面方程为\pickout{A}
\options
{$x-z=-1$}
{$x+z=1$}
{$x-z=1$}
{$x+y-z=1-\frac{\pi}{2}$}
\end{problem}
\begin{note}
$z_x=\mathrm{e}^x\sin y=1$，$z_y=\mathrm{e}^x\cos y=0$。切平面：$z-1=1\cdot(x-0)+0\cdot\bigl(y-\frac{\pi}{2}\bigr)$，即 $x-z=-1$。
\end{note}
\bigskip

\begin{problem}
曲面 $x^2+y^2=\mathrm{e}^{z}$ 在点 $(1,0,0)$ 处的一个法向量为\pickout{B}
\options
{$(1,0,-1)$}
{$(2,0,-1)$}
{$(2,0,1)$}
{$(1,0,-\mathrm{e})$}
\end{problem}
\begin{note}
设 $F=x^2+y^2-\mathrm{e}^{z}$，则 $\nabla F=(2x,\,2y,\,-\mathrm{e}^{z})\big|_{(1,0,0)}=(2,0,-1)$。
\end{note}
\bigskip

\begin{problem}
函数 $f(x,y)=\mathrm{e}^{x^2}+\mathrm{e}^{y^2}$ 在点 $(0,0)$ 处\pickout{B}
\options
{取极大值}
{取极小值}
{是鞍点}
{不是驻点}
\end{problem}
\begin{note}
$f_x=2x\mathrm{e}^{x^2}$，$f_y=2y\mathrm{e}^{y^2}$，在 $(0,0)$ 处均为 $0$。
$f_{xx}=2\mathrm{e}^{x^2}+4x^2\mathrm{e}^{x^2}=2$，$f_{yy}=2$，$f_{xy}=0$。
$AC-B^2=4>0$ 且 $A=2>0$，故为极小值点。
\end{note}
\bigskip

\begin{problem}
函数 $f(x,y)=x^2+\cos y$ 在点 $(0,0)$ 处\pickout{C}
\options
{取极大值}
{取极小值}
{是鞍点}
{不是驻点}
\end{problem}
\begin{note}
$f_x=2x=0$，$f_y=-\sin y=0$，$(0,0)$ 是驻点。
$f_{xx}=2$，$f_{yy}=-\cos y=-1$，$f_{xy}=0$。
海瑟矩阵 $H=\begin{pmatrix}2&0\\0&-1\end{pmatrix}$ 不定（特征值一正一负），故为鞍点。
\end{note}
\bigskip

\begin{problem}
曲线 $\begin{cases}z=\mathrm{e}^{x+y},\\x^2+y^2=2\end{cases}$ 在点 $(1,1,\mathrm{e}^{2})$ 处的切向量可取为\pickout{A}
\options
{$(1,-1,0)$}
{$(1,1,0)$}
{$(1,1,\mathrm{e}^2)$}
{$(1,0,-\mathrm{e}^2)$}
\end{problem}
\begin{note}
设 $F=\mathrm{e}^{x+y}-z$，$G=x^2+y^2-2$，则
$\nabla F=(\mathrm{e}^{x+y},\,\mathrm{e}^{x+y},\,-1)\big|_{(1,1,\mathrm{e}^2)}=(\mathrm{e}^2,\,\mathrm{e}^2,\,-1)$，
$\nabla G=(2x,\,2y,\,0)\big|_{(1,1,\mathrm{e}^2)}=(2,\,2,\,0)$。
切向量 $s=\nabla F\times\nabla G=(2\mathrm{e}^2,\,-2\mathrm{e}^2,\,0)=2\mathrm{e}^2(1,-1,0)$。
\end{note}
\bigskip

% =================================================
%       解答题：极值应用
% =================================================
\newpage
\makepart{解答题}{每小题~11~分，共~22~分}

\begin{problem}
要制作一个容积为 $32\,\mathrm{m}^3$ 的无盖长方体容器，底面为正方形。求使所用材料最省的底面边长和高，并求此时的最小表面积。
\end{problem}
\begin{note}
设底面边长为 $x$，高为 $h$。目标函数 $S=x^2+4xh$（无盖，仅有底面与四个侧面），约束条件 $x^2h=32$。使用拉格朗日乘数法求解。
\end{note}
\begin{solution}
设底面边长为 $x$（$x>0$），高为 $h$（$h>0$），则约束条件为
\[
x^2h=32.
\]
无盖容器的表面积（即用料面积）为
\[
S=x^2+4xh.
\]
构造拉格朗日函数
\[
\mathcal{L}(x,h,\lambda)=x^2+4xh+\lambda\bigl(x^2h-32\bigr).
\]
求偏导并令其为零：
\begin{align*}
\frac{\partial\mathcal{L}}{\partial x}&=2x+4h+2\lambda xh=0,\tag{1}\\
\frac{\partial\mathcal{L}}{\partial h}&=4x+\lambda x^2=0,\tag{2}\\
\frac{\partial\mathcal{L}}{\partial\lambda}&=x^2h-32=0.\tag{3}
\end{align*}
由式 $(2)$，因 $x>0$，可得 $\lambda=-\dfrac{4}{x}$。
将其代入式 $(1)$：
\[
2x+4h+2\left(-\frac{4}{x}\right)xh=2x+4h-8h=2x-4h=0,
\]
故 $x=2h$。

将 $x=2h$ 代入约束式 $(3)$：
\[
(2h)^2\cdot h=4h^3=32\quad\Longrightarrow\quad h^3=8,
\]
得 $h=2$，进而 $x=4$。

因此，当底面边长为 $\boxed{4\,\mathrm{m}}$、高为 $\boxed{2\,\mathrm{m}}$ 时，所用材料最省。
此时的最小表面积为
\[
S_{\min}=4^2+4\times4\times2=16+32=\boxed{48\,\mathrm{m}^2}.
\]
\end{solution}

\vspace{1cm}

\begin{problem}
某公司在两个市场销售同一种服务，售价分别为 $p_1$ 和 $p_2$（单位：百元）。两个市场的需求量分别为 $q_1=24-2p_1$ 和 $q_2=20-p_2$。该服务的固定成本为 $10$ 百元，每单位可变成本为 $4$ 百元。由于市场调控政策，两个市场的售价必须满足 $p_1-p_2=2$。求使公司利润最大的售价 $p_1,p_2$，并求最大利润。
\end{problem}
\begin{note}
先写出利润函数，再以 $p_1-p_2=2$ 为约束，用拉格朗日乘数法求解。也可将 $p_1=p_2+2$ 代入化为无条件极值求解。
\end{note}
\begin{solution}
总成本为 $C=4(q_1+q_2)+10$，收入为 $R=p_1q_1+p_2q_2$，故利润
\begin{align*}
L&=p_1(24-2p_1)+p_2(20-p_2)-\bigl[4(24-2p_1+20-p_2)+10\bigr]\\
&=-2p_1^2-p_2^2+32p_1+24p_2-186.
\end{align*}
约束条件为 $g(p_1,p_2)=p_1-p_2-2=0$。

构造拉格朗日函数
\[
\mathcal{L}(p_1,p_2,\lambda)=-2p_1^2-p_2^2+32p_1+24p_2-186+\lambda(p_1-p_2-2).
\]
求偏导并令其为零：
\begin{align*}
\frac{\partial\mathcal{L}}{\partial p_1}&=-4p_1+32+\lambda=0,\tag{1}\\
\frac{\partial\mathcal{L}}{\partial p_2}&=-2p_2+24-\lambda=0,\tag{2}\\
\frac{\partial\mathcal{L}}{\partial\lambda}&=p_1-p_2-2=0.\tag{3}
\end{align*}

将式 $(1)$ 与式 $(2)$ 相加，消去 $\lambda$：
\[
-4p_1-2p_2+56=0\quad\Longrightarrow\quad 2p_1+p_2=28.
\]
结合式 $(3)$ 的 $p_1=p_2+2$，代入上式得
\[
2(p_2+2)+p_2=28\quad\Longrightarrow\quad 3p_2=24,
\]
解得 $p_2=8$，进而 $p_1=10$。

验证：此时 $q_1=24-20=4$，$q_2=20-8=12$。
收入 $R=10\times4+8\times12=136$（百元），
成本 $C=4\times(4+12)+10=74$（百元），
最大利润为
\[
L_{\max}=136-74=\boxed{62\text{（百元）}}.
\]

故最优售价为 $\boxed{p_1=10}$ 百元、$\boxed{p_2=8}$ 百元，最大利润为 $\boxed{62}$ 百元。
\end{solution}

\end{document}
